%------------------------------------------------------------------------------%
\section{Notação Matemática}
\label{sec:Notacao_Matematica}

%------------------------------------------------------------------------------%
\subsection*{Símbolos Matemáticos}

A lista a seguir inclui vários símbolos matemáticos comumente utilizados e seus
significados

{\centering
\begin{tabular}{lp{15cm}}
  \emph{$\infty\;$} &
  Indica uma quantidade infinita. Observe que $\infty$ não é um número e portanto
  não podemos realizar operações algébricas com ele.
  \\
  \emph{$\forall\;$} &
  Para todo
  \\
  \emph{$\exists\;$} &
  Existe
  \\
  \emph{$\to\;$} &
  Esse símbolo pode indicar um limite, $x\to 0$, ou ser usado na definição de uma função
  $f:\R\to\R$. Ele \emph{não} pode substituir outros símbolos como a igualdade e também
  \emph{não} deve ser usado para indicar o sentido de leitura.
\end{tabular}\par}

A seguir apresentamos símbolos que indicam relações. Note que cada um tem um
significado específico, um \emph{não} pode substituir o outro.

{\centering
\begin{tabular}{lp{15cm}}
  \emph{$=\;$} &
  Símbolo de igualdade, ele indica que as duas expressões são completamente equivalentes.
  \\
  \emph{$\Leftrightarrow\;$} &
  Esse símbolo significa ``se, e somente se''.
  Ele só pode ser utilizado entre duas afirmações lógicas, isso é, que podem
  ser verdadeiras ou falsas, e indica que ambas são simultaneamente
  verdadeiras ou simultaneamente falsas.
  \\
  \emph{$\Rightarrow\;$} &
  Assim como $\Leftrightarrow$, esse símbolo indica uma relação lógica entre duas
  afirmações. Porém, seu significado é mais sutil, ele indica que se a primeira
  afirmação for verdadeira a segunda também precisa ser. Se a primeira afirmação
  for falsa não sabemos nada sobre a segunda.
  \\
  \emph{$\Leftarrow\;$} &
  Tem o mesmo significado que $\Rightarrow$, porém com os papeis trocados para as afirmações.
\end{tabular}\par}

%------------------------------------------------------------------------------%
\subsection*{Conjuntos e Intervalos}
\label{sub:conjuntos_intervalos}

Os conjuntos numéricos mais utilizados são:

{\centering
\begin{tabular}{llp{10cm}}
  \emph{$\N\;$} & Números naturais  & $\{1, 2, 3, \dots\}$ \\
  \emph{$\Z\;$} & Números inteiros  & $\{\dots, -3, -2, -1, 0, 1, 2, 3, \dots\}$ \\
  \emph{$\Q\;$} & Números racionais & Razões, ou frações, de números inteiros \\
  \emph{$\R\;$} & Números reais     & A caracterização dos reais é mais complexa,
                                      para o Cálculo, podemos pensar que eles correspondem
                                      aos pontos da reta e que não há buracos entre eles \\
  \emph{$\C\;$} & Números complexos & Extensão dos reais para incluir raízes de números negativos
\end{tabular}\par}

As operações com conjuntos mais comuns são

{\centering
\begin{tabular}{lp{15cm}}
  \emph{$\in\;$}     & Indica que um elemento pertence a um conjunto \\
  \emph{$\cup\;$}    & União de conjuntos, representa os elementos que pertencem a um ou outro conjunto \\
  \emph{$\cap\;$}    & Intersecção de conjuntos, representa os elementos que pertencem a ambos conjuntos \\
  \emph{$\subset\;$} & Subconjunto, indica que todos os elementos do conjunto da esquerda pertencem
                       também ao conjunto da direita.
\end{tabular}\par}

\emph{Intervalos}\index{Intervalo}
são um caso particular de subconjuntos dos números reais,
eles indicam um conjunto de números reais ``consecutivos'' e ``sem buracos''.
Normalmente os intervalos são definidos por desigualdades
\[
\begin{array}{cl}
  I & = \left\{\, x\in\R \st 3 < x < 5        \,\right\} \\
  J & = \left\{\, x\in\R \st -1 \leq x \leq 2 \,\right\} \\
  K & = \left\{\, x\in\R \st \abs{x} < 2      \,\right\}
\end{array}
\]
Como esses conjuntos são muito utilizados no Cálculo existe uma notação
específica para eles. A notação para intervalos precisa indicar quais
são seus extremos e se esses extremos pertencem ou não ao intervalo,
um colchete $[$ ou $]$ indica que o extremo correspondente pertence
ao intervalo enquanto que um parenteses $($ ou $)$ indica que os extremos
não pertence ao intervalo
\[
\begin{array}{cl}
  (3, 5)  & = \left\{\, x\in\R \st 3 < x < 5        \,\right\} \\[1pt]
  [-1, 2] & = \left\{\, x\in\R \st -1 \leq x \leq 2 \,\right\} \\[1pt]
  [-2, 2] & = \left\{\, x\in\R \st \abs{x} < 2      \,\right\} \\[1pt]
  (a, b]  & = \left\{\, x\in\R \st a <    x \leq b  \,\right\} \\[1pt]
  [a, b)  & = \left\{\, x\in\R \st a \leq x <    b  \,\right\}
\end{array}
\]
Quando ambos extremos pertencem ao intervalo dizemos que ele é um
\emph{Intervalo Fechado}\index{Intervalo!fechado} quando
nenhum extremo pertence ao intervalo dizemos que ele é um
\emph{Intervalo Aberto}\index{Intervalo!aberto}.

Um intervalo aberto que contém um ponto $a$ é chamado de
\emph{Vizinhança}\index{Vizinhança} de $a$.

%------------------------------------------------------------------------------%

\begin{definicao}{Ínfimo e Supremo}{def:inf-sup}
  O \emph{Ínfimo}\index{Infimo@Ínfimo} de $A\subset\R$, $\inf(A)$,
  é o maior número real menor ou igual a todos os elementos de $A$.

  O \emph{Supremo}\index{Supremo} de $A\subset\R$, $\sup(A)$,
  que é o menor número real maior ou igual a todos os elementos de $A$.
\end{definicao}

Para observamos a diferença entre o ínfimo e o mínimo considere
o intervalo aberto entre zero e um, $A=(0, 1)$.
Nenhum elemento $x\in A$ pode ser o mínimo de $A$, pois $\sfrac{x}{2}$
é menor do que $x$ e pertence a $A$, portanto, esse conjunto não possui
um menor elemento. O zero seria o candidato para ser o mínimo, mas
ele não pertence a $A$, nesse caso dizemos que zero é o ínfimo de $A$.
Uma análise equivalente vai nos mostrar que $A$ também não possui um
máximo e que 1 é seu supremo. Podemos escrever em tão que
\begin{align*}
  \inf\big((0, 1)\big) & = 0  &  \min\big((0, 1)\big) \quad \text{não existe} \\
  \sup\big((0, 1)\big) & = 1  &  \max\big((0, 1)\big) \quad \text{não existe}
\end{align*}
Em um intervalo fechado o ínfimo coincide com o mínimo e o supremo
coincide com o máximo, por exemplo, considere $[0, 1]$
\begin{align*}
  \inf\big([0, 1]\big) & = \,\min\big([0, 1]\big) = 0 \\
  \sup\big([0, 1]\big) & =   \max\big([0, 1]\big) = 1
\end{align*}

%------------------------------------------------------------------------------%
